\documentclass[a4paper,11pt]{article}
\usepackage{jcappub}
\usepackage{lineno}
\linenumbers

\title{\boldmath A focused CEvNS and neutrino-electron benchmark for a 200 m$^3$ CF$_4$ gas TPC}

\author{D. TODO}
\affiliation{TODO institution,\\
TODO address, Country}

\emailAdd{TODO@example.org}

\abstract{
We present a focused feasibility benchmark for solar and core-collapse
supernova neutrino detection in a CYGNUS-K-like gas time projection chamber.
The benchmark fixes a 10 m $\times$ 10 m $\times$ 2 m fiducial volume
(200 m$^3$) filled with CF$_4$ at the existing 10 atm gas-table point, used as
the first proxy for a 10 bar configuration. We treat neutrino-electron
scattering and coherent elastic neutrino-nucleus scattering (CEvNS) separately:
directional electron-recoil information is used only for the neutrino-electron
channel, while CEvNS is assumed to be non-directional and threshold-limited.
For the CEvNS baseline we use a 1.0 keVee visible threshold converted to
isotope-dependent true nuclear-recoil thresholds with TRIM quenching curves for
C and F in CF$_4$. This paper is not a complete background-limited sensitivity
study; it defines the benchmark, reports first-pass rates and burst counts, and
identifies the detector-response ingredients needed for a full feasibility
assessment.
}

\begin{document}
\maketitle
\flushbottom

\section{Introduction}
\label{sec:introduction}

Directional gas time projection chambers are usually discussed in the context
of rare-event searches, but the same low-threshold gaseous target can also
provide a useful neutrino benchmark. The purpose of this note is deliberately
narrow: we define one fixed CYGNUS-K-like configuration and ask what solar and
supernova neutrino signals are expected before a full background model is
introduced.

The key detector assumption is channel-dependent. For neutrino-electron
scattering, electron range and diffusion determine the accepted electron-recoil
window, and angular information can improve neutrino-energy reconstruction. For
CEvNS, however, no event-by-event directionality is assumed. The CEvNS
benchmark is therefore controlled by a visible-energy threshold and by nuclear
quenching, not by the electron-recoil directional threshold.

The result should be read as a compact feasibility benchmark rather than a
discovery-reach paper. Backgrounds, efficiency turn-on curves, trigger
implementation, and systematic uncertainty propagation are listed explicitly as
future work in section~\ref{sec:limitations}.

\section{Benchmark detector and target}
\label{sec:benchmark}

The baseline detector is the rectangular fiducial geometry used by the analysis
repository:
\begin{equation}
L \times W \times H = 2~{\rm m} \times 10~{\rm m} \times 10~{\rm m}
= 200~{\rm m}^3 \, .
\end{equation}
The target gas is CF$_4$ at the 10 atm entry in the gas-density table. In this
draft we refer to this as the 10 bar benchmark for readability, but the final
paper must either convert the gas table to exactly 10 bar or consistently use
10 atm throughout.

\begin{table}[htbp]
\centering
\begin{tabular}{ll}
\hline
Quantity & Baseline value \\
\hline
Fiducial volume & 200 m$^3$ \\
Geometry & 10 m $\times$ 10 m $\times$ 2 m \\
Gas & CF$_4$ \\
Gas table row & 10 atm (10 bar proxy) \\
CEvNS threshold mode & 1.0 keVee converted with TRIM QF \\
CEvNS directionality & none assumed \\
$\nu$-e directionality & allowed for electron recoils only \\
Supernova reference model & SN1987A-like fluence at 10 kpc \\
\hline
\end{tabular}
\caption{Fixed run card for the first focused benchmark.}
\label{tab:run-card}
\end{table}

\section{Directionality and threshold convention}
\label{sec:thresholds}

For the neutrino-electron channel, the accepted electron recoil energy window is
derived from electron range and diffusion. This is a reconstruction and
topology requirement on electron tracks.

For CEvNS, the observable is instead a quenched nuclear recoil. The baseline
visible threshold is
\begin{equation}
E_{\rm ee}^{\rm thr} = 1.0~{\rm keVee} \, .
\end{equation}
For isotope $i$, the corresponding true nuclear recoil threshold is obtained by
solving
\begin{equation}
Q_i(T_N) T_N = E_{\rm ee}^{\rm thr} \, ,
\label{eq:qf-threshold}
\end{equation}
where $Q_i$ is the TRIM quenching factor for C or F recoils in CF$_4$. The
current TRIM tables imply approximately 3.0 keVnr for carbon and 3.6 keVnr for
fluorine at 1.0 keVee. The analysis code forbids extrapolation outside the TRIM
tables for the baseline.

\begin{figure}[htbp]
\centering
\fbox{\begin{minipage}{0.88\textwidth}
TODO: replace this box with the quenching-factor plot generated from
\texttt{quenching\_factor/C\_in\_CF4.csv} and
\texttt{quenching\_factor/F\_in\_CF4.csv}. Mark the 1.0 keVee crossing for C
and F.
\end{minipage}}
\caption{TRIM quenching curves used to convert the CEvNS visible threshold to
isotope-dependent true nuclear-recoil thresholds.}
\label{fig:qf}
\end{figure}

\section{Neutrino inputs and interaction model}
\label{sec:model}

The solar calculation uses the repository solar flux tables and the compact
flavor split described in the code documentation. The supernova calculation
uses the analytic time-integrated fluence model currently labelled
SN1987A-like, with the reference fluence generated at 10 kpc.

The neutrino-electron channel uses Standard Model elastic scattering on target
electrons. For CEvNS we use the vector weak charge for C and F, with the
implemented axial contribution for $^{19}$F. The current hydrogen axial
limitation is irrelevant for the CF$_4$ benchmark.

\section{Focused benchmark results}
\label{sec:results}

This section records the regenerated CF$_4$ 10 atm outputs from the run card in
appendix~\ref{app:run-card}. These values should be refreshed whenever the
detector response, flux inputs, or threshold convention changes.

\begin{table}[htbp]
\centering
\begin{tabular}{lr}
\hline
Observable & Value \\
\hline
Solar $\nu$-e total rate & $5.62 \times 10^3$ yr$^{-1}$ \\
Solar CEvNS total rate & $2.59 \times 10^2$ yr$^{-1}$ \\
SN $\nu$-e counts at 10 kpc & $6.09 \times 10^{-2}$ \\
SN CEvNS counts at 10 kpc & $17.1$ \\
Distance where CEvNS gives 1 event & $41.4$ kpc \\
\hline
\end{tabular}
\caption{Headline benchmark quantities for CF$_4$ at 10 atm in 200 m$^3$,
using the 1.0 keVee QF-converted CEvNS threshold.}
\label{tab:headline}
\end{table}

\begin{figure}[htbp]
\centering
\fbox{\begin{minipage}{0.88\textwidth}
TODO: insert the solar CEvNS recoil and accepted-neutrino-energy spectra from
\texttt{solar\_nu\_gas\_target\_rates/cevns/cf4\_10atm/}.
\end{minipage}}
\caption{Solar CEvNS spectra for the focused CF$_4$ benchmark.}
\label{fig:solar-cevns}
\end{figure}

\begin{figure}[htbp]
\centering
\fbox{\begin{minipage}{0.88\textwidth}
TODO: insert the supernova counts-versus-distance plot from
\texttt{supernova\_distance\_gas\_scan/SN1987A\_like/cf4\_10atm/}.
\end{minipage}}
\caption{Expected supernova event counts versus distance for CF$_4$ at 10 atm.}
\label{fig:sn-distance}
\end{figure}

\section{Limitations and next steps}
\label{sec:limitations}

The focused benchmark intentionally omits several ingredients that are required
for a complete feasibility paper:
\begin{itemize}
\item detector efficiency and trigger turn-on near the CEvNS threshold;
\item background rates, including radiogenic, cosmogenic, and instrumental
backgrounds;
\item a validated nuclear-recoil energy-resolution model;
\item uncertainty propagation for solar fluxes, supernova spectra, form
factors, axial response, quenching, gas density, and threshold;
\item conversion from the 10 atm proxy to an exactly specified 10 bar operating
point.
\end{itemize}

The next paper iteration should convert this benchmark into a sensitivity
study by adding efficiency and background models and by reporting detection
probabilities or confidence intervals instead of only raw rates and counts.

\section{Conclusions}
\label{sec:conclusions}

We have defined a focused, reproducible benchmark for a 200 m$^3$ CF$_4$
CYGNUS-K-like detector. The central methodological point is that CEvNS is
treated as non-directional and therefore uses a quenching-converted visible
threshold, while directionality remains part of the neutrino-electron channel.
Once the regenerated CF$_4$ 10 atm results are inserted, this draft can serve
as the seed for a compact JCAP feasibility paper.

\appendix

\section{Run card}
\label{app:run-card}

The benchmark outputs should be regenerated with a single consistent threshold
configuration:
\begin{verbatim}
python3 plotFluxes.py
python3 gasTargetRates.py --gas CF4 --pressure-atm 10
python3 recoNuEnergyComparison.py --gas CF4 --pressure-atm 10

python3 generate_sn_fluence.py
python3 supernovaGasTargetEvents.py --gas CF4 --pressure-atm 10
python3 supernovaRecoNuEnergyComparison.py --gas CF4 --pressure-atm 10

python3 supernovaDistanceGasScan.py \
  --gas CF4 \
  --pressure-atm 10 \
  --cevns-threshold-mode qf_electron_equivalent \
  --cevns-visible-threshold-kevee 1.0
\end{verbatim}

The CEvNS configuration block used for the paper should contain
\begin{verbatim}
"threshold_mode": "qf_electron_equivalent",
"visible_threshold_keVee": 1.0,
"quenching_dir": "quenching_factor"
\end{verbatim}
with CEvNS enabled.

\section{CEvNS threshold conversion}
\label{app:qf}

The TRIM quenching files are two-column CSV files. The first column is true
nuclear recoil energy in keVnr and the second column is
$Q(T_N)=E_{\rm ee}/T_N$. The code sorts each curve by true recoil energy,
checks that $Q(T_N)T_N$ is strictly increasing, and then linearly interpolates
the solution to equation~\eqref{eq:qf-threshold}. Extrapolation is not allowed
for the baseline.

\acknowledgments

TODO.

\begin{thebibliography}{99}

\bibitem{Freedman1974}
D.Z. Freedman,
\emph{Coherent effects of a weak neutral current},
Phys. Rev. D {\bf 9} (1974) 1389.

\bibitem{COHERENT2017}
COHERENT collaboration,
\emph{Observation of coherent elastic neutrino-nucleus scattering},
Science {\bf 357} (2017) 1123.

\bibitem{Angloher2025}
G. Angloher et al.,
\emph{Neutrino flux sensitivity to the next galactic core-collapse supernova in COSINUS},
JCAP {\bf 03} (2025) 037.

\bibitem{CYGNUS}
CYGNUS collaboration,
\emph{CYGNUS: Feasibility of a nuclear recoil observatory with directional sensitivity to dark matter and neutrinos},
TODO update bibliographic data before submission.

\end{thebibliography}

\end{document}
